% =============================================================================
% Title Block  (spans both columns via \twocolumn[...])
% =============================================================================
\twocolumn[{
    \begin{flushleft}
        {\small
        Department of Software Engineering\\
        Faculty of Sciences and Techniques Tangier\\
        Abdelmalek Essadi University, Morocco\\
        19/02/2026\par
        }

        \vspace{2.5em}

        {\color{titleblue}\Huge\bfseries Automatic Detection of Water Losses in  \\[0.2em]
        Urban Distribution and Irrigation Networks\\[0.2em]
        Using Deep Learning\\[0.2em]
        \par}

        \vspace{1.5em}

        {\large\bfseries Supervised by: Pr. Mostafa EZZIYANI\\[0.8em]
        Made By: ZARQI Ezzoubair, El Gharib Mahmoud\par}
    \end{flushleft}

    \vspace{1.5em}

    % Abstract
    \begin{leftbar}
        \textbf{\textcolor{headingblue}{ABSTRACT}} Water losses due to leaks in urban
        distribution and irrigation networks represent a critical global challenge, leading
        to significant economic losses and resource waste. Traditional leak detection methods
        based on manual inspection and threshold-based rules are often insufficient, as leaks
        produce subtle, sustained deviations in sensor readings rather than easily detectable
        sharp anomalies. In this study, we propose a deep learning approach for automatic
        water leak detection using multivariate time-series data collected from pressure, flow
        rate, and temperature sensors. We transform the raw sensor data into overlapping
        sliding windows and train a one-dimensional Convolutional Neural Network (1D-CNN)
        with batch normalization and global average pooling. To ensure an unbiased evaluation,
        we employ \textbf{sensor-wise GroupKFold cross-validation} that guarantees zero data
        leakage between training and validation sets. To address class imbalance, we apply
        balanced class weighting. Evaluated across 5 folds with strict sensor-wise separation,
        our model achieves 71.30\% accuracy, 74.07\% precision, 85.71\% recall, and a 0.79
        F1-score, providing a realistic baseline for leak detection on a small-scale dataset
        of 1{,}000 observations.
    \end{leftbar}

    \vspace{1em}

    % Keywords
    \begin{leftbar}
        \textbf{\textcolor{headingblue}{KEY-WORDS}} Water Leak Detection, Deep Learning,
        1D Convolutional Neural Network, Time-Series Classification, Sliding Window,
        Sensor Data, Urban Water Distribution, Class Imbalance;
    \end{leftbar}

    \vspace{2.5em}
}]
